%!TEX root = TheoreticalPhysicsC.tex

\section{Lecture: Particle in a Box}
A single particle all alone in the universe models no real system. A particle in a box of finite size, by contrast, is a serviceable effective model: for an electron in a quantum dot, say, or for a photon in a cavity resonator. Along the way we see here for the first time where the word ``quantum'' comes from.

If the freedom of motion of a particle is restricted by external walls, this shows up first of all as a restriction on the possible values of the position variable $x$. If the particle can move only inside a region $G \subset \mathbb{R}^3$, we first modify the Hilbert space to $\mathcal{H}\equiv L^2(G)$, i.e.\ wave functions are defined only on $G$, and in the
inner product the integration 
\begin{equation}
	\langle \Phi|\Psi\rangle \equiv \iiint_G \Phi^*(x,y,z)\Psi(x,y,z)\, dxdydz
\end{equation}
is likewise restricted to $G$. 

For the Hamiltonian we again try the ansatz 
\begin{equation}
	H = -\frac{\hbar^2}{2m}\nabla^2
\end{equation}
This, however, goes wrong ($H$ will not be Hermitian). One can restore the Hermiticity of $H$ by restricting the domain of $H$ to those functions which satisfy certain \emph{boundary conditions}. These boundary conditions are ultimately an idealised description of the interaction between particle and wall. A frequently used boundary condition for the Hamiltonian is the so-called \emph{Dirichlet boundary condition}: only functions which vanish on the boundary of $G$ are admitted into the domain of $H$.

How crude this idealisation is becomes clear if one tries to describe the wall honestly: it consists of an astronomical number of atoms, every single one of which would have to be modelled as well -- that is surface physics. Strictly speaking the boundary condition is even wrong: a tiny probability of finding the particle outside $G$ always remains, in the gaps between the atoms, for instance. The procedure is a familiar one, though: in electrodynamics too, certain field components vanish at conductor surfaces.

In order to solve the dynamics of $H$ it is extremely helpful to compute the eigenvalues and the associated eigenvectors of $H$. (Then $U_t=e^{-\frac{it}{\hbar}\widehat{H}}$ is given by nothing but a multiplication operator.) Our old eigenvectors are of no further use for this: a plane wave is non-zero everywhere and cannot possibly satisfy the boundary condition. What worked for the free particle fails here as a matter of principle. We begin by computing the eigenfunctions for the simplest case: a ``one-dimensional'' particle in the interval $[0,L] \subset \mathbb{R}$. That is, we first compute the solutions of the equation
\begin{equation}
	\widehat{H}|\Psi\rangle = E|\Psi\rangle
\end{equation}
for all possible eigenvalues $E$. In the literature it is usually called the \emph{stationary Schrödinger equation} -- really an unnecessary name for a plain eigenvector equation. To get any further we shall rewrite this equation as an equation for the wave function
\begin{equation}
	\langle x|\widehat{H}|\Psi\rangle = -\frac{\hbar^2}{2m}\frac{d^2}{dx^2}\Psi(x) = E\Psi(x) = E\langle x|\Psi\rangle.
\end{equation}
Without further restrictions the Hamiltonian is not Hermitian. We must also impose the boundary condition 
\begin{equation}
	\Psi(0) = \Psi(L) = 0
\end{equation} 
This is a linear differential equation with constant coefficients, and is therefore
solved by the ansatz
\begin{equation}
	\Psi(x) = e^{\mu x}, \quad\text{with}\quad \mu\in \mathbb{C}.
\end{equation}
Substituting gives the equation
\begin{equation}
	-\frac{\hbar^2}{2m} \mu^2 e^{\mu x} = E e^{\mu x},
\end{equation}
so
\begin{equation}
	\mu = \pm i \sqrt{\frac{2mE}{\hbar^2}}, 
\end{equation}
and hence two linearly independent solutions.

The general solution of the differential equation is therefore: 
\begin{equation}
	\Psi(x) = A e^{\mu x} + B e^{-\mu x},
\end{equation}
where $\mu$ is one of the two solutions, and $A,B \in \mathbb{C}$. The boundary condition $\Psi(0)= 0$ forces $B = -A$. So we are left with the equation 
\begin{equation}
	\Psi(L) =A e^{\mu L} + B e^{-\mu L} = 0,
\end{equation}
or equivalently: 
\begin{equation}
	e^{2\mu L} = 1,
\end{equation}
from which the possible values of $\mu$ (and with them of $E$) follow. Real $\mu$ are ruled out at once,
so $E \ge 0$. For imaginary $\mu$ there are solutions, namely 
\begin{equation}
	2\mu L = 2\pi i n,
\end{equation}
where $n \in \mathbb{Z}$. Thus $\mu = n\pi i/L$, and the eigenfunction $\Psi(x) = A \sin(n\pi x/L)$.
Clearly $n$ and $-n$ give the same eigenfunction, and $n = 0$ is excluded (there $\Psi'' =0$, so that $
\Psi(x) = ax+b$ is linear, which does not sit well with the boundary conditions). After normalisation we therefore obtain the eigenfunctions and eigenvalues
\begin{equation}
	\phi_n(x) = \sqrt{\frac{2}{L}}\sin\left(\frac{n\pi}{L} x\right)\quad \text{and}\quad E_n = \frac{\hbar^2 \pi^2}{2mL^2}n^2\quad \text{for $n=1,2,\ldots$.}
\end{equation}
Pictorially, $\phi_1$ is half a sine wave across the interval, $\phi_2$ a whole one, and so on; you had best check the normalisation integral for yourself.

\begin{merksatz}[Where the name ``quantum mechanics'' comes from]
	For a free particle every energy was allowed. In the box, all that survives of this continuum is the discrete sequence $E_n\propto n^2$: restricting the position leads as a rule to the \emph{quantisation} of the energy eigenvalues -- hence the name of the whole theory.
\end{merksatz}

% fig_vl11_kasten.tex — Teilchen im Kasten: Eigenfunktionen auf ihren
% Energieniveaus, und die quadratische Aufspreizung des Spektrums.
%
% Reference figure for the notes: the other lecture figures follow the
% conventions established here.
%
%   * geometry and curves come from the closed-form solution derived in
%     the lecture, never from a freehand sketch:
%         psi_n(x) = sqrt(2/L) sin(n pi x / L),
%         E_n      = n^2 pi^2 hbar^2 / (2 m L^2),
%     so the node count, the symmetry and the level spacing are forced
%     to be right rather than drawn by eye;
%   * each eigenfunction is drawn about its own energy — the standard way
%     of showing that psi_n belongs to E_n;
%   * colour encodes the quantum number n (the ordered ramp of
%     `tpclevels`), and nothing else.  The energy levels themselves are
%     neutral grey rules: they are reference lines, not data;
%   * the classically forbidden region is hatched, not filled, so it
%     survives greyscale printing.
%
% Panel (a) the first four eigenfunctions in the well.
% Panel (b) the spectrum E_n/E_1, with the oscillator's evenly spaced
%           ladder alongside it — the contrast between a quadratic and a
%           linear spectrum is the point students should carry away.

\begin{center}
\begin{tikzpicture}

% ─────────────────────────────────────────────────────────────
%  (a)  Eigenfunctions drawn about their levels
% ─────────────────────────────────────────────────────────────
\begin{axis}[
    name=wells,
    tpcplot,
    width=8.2cm, height=6.8cm,
    axis lines=none,
    xmin=-0.34, xmax=1.34,
    ymin=-1.2,  ymax=19.4,
    clip=false,
    title={(a) \tpcifenglish{$\psi_n$ on their levels}%
                       {$\psi_n$ auf ihren Niveaus}},
    domain=0:1, samples=201,
]

  % Classically forbidden regions, x<0 and x>L, where V = infinity
  \addplot[forbidden, draw=none, domain=-0.34:0, samples=2] {19.4} \closedcycle;
  \addplot[forbidden, draw=none, domain=1:1.34,  samples=2] {19.4} \closedcycle;

  % The well itself
  \draw[potential] (axis cs:0,-1.2) -- (axis cs:0,19.4);
  \draw[potential] (axis cs:1,-1.2) -- (axis cs:1,19.4);
  \draw[potential] (axis cs:0,0)    -- (axis cs:1,0);

  % Energy levels E_n = n^2 E_1: neutral reference rules
  \addplot[levellinegrey, domain=0:1, samples=2] {1};
  \addplot[levellinegrey, domain=0:1, samples=2] {4};
  \addplot[levellinegrey, domain=0:1, samples=2] {9};
  \addplot[levellinegrey, domain=0:1, samples=2] {16};

  % Eigenfunctions about those levels.  Amplitude 1.3 is the largest that
  % keeps psi_2 clear of E_1, the tightest gap (E_2 - E_1 = 3 E_1).
  \addplot[wavefn, spacecadet]   {1  + 1.3*sin(deg(1*pi*x))};
  \addplot[wavefn, cgblue]       {4  + 1.3*sin(deg(2*pi*x))};
  \addplot[wavefn, munsell]      {9  + 1.3*sin(deg(3*pi*x))};
  \addplot[wavefn, birrenorange] {16 + 1.3*sin(deg(4*pi*x))};

  % Level labels.  Written out rather than looped: \foreach defers
  % expansion inside a pgfplots axis.
  \node[lbl, anchor=east] at (axis cs:-0.06,1)  {$E_{1}$};
  \node[lbl, anchor=east] at (axis cs:-0.06,4)  {$E_{2}$};
  \node[lbl, anchor=east] at (axis cs:-0.06,9)  {$E_{3}$};
  \node[lbl, anchor=east] at (axis cs:-0.06,16) {$E_{4}$};

  % Node counts: psi_n has n-1 interior nodes
  \node[lbl, text=birrengray, font=\scriptsize, anchor=west]
    at (axis cs:1.03,1)  {$n=1$};
  \node[lbl, text=birrengray, font=\scriptsize, anchor=west]
    at (axis cs:1.03,4)  {$n=2$};
  \node[lbl, text=birrengray, font=\scriptsize, anchor=west]
    at (axis cs:1.03,9)  {$n=3$};
  \node[lbl, text=birrengray, font=\scriptsize, anchor=west]
    at (axis cs:1.03,16) {$n=4$};

  % Axis furniture
  \node[lbl, anchor=north] at (axis cs:0,-0.3)    {$0$};
  \node[lbl, anchor=north] at (axis cs:1,-0.3)    {$L$};
  \node[lbl, anchor=north] at (axis cs:0.5,-0.62) {$x$};
  \node[lbl, text=birrengray, font=\scriptsize, anchor=south]
    at (axis cs:-0.17,17.6) {$V=\infty$};
\end{axis}

% ─────────────────────────────────────────────────────────────
%  (b)  Spectrum: the n^2 ladder against the oscillator's linear one
% ─────────────────────────────────────────────────────────────
\begin{axis}[
    tpcplot,
    at={(wells.south east)}, anchor=south west, xshift=17mm,
    width=6.2cm, height=6.8cm,
    xlabel={$n$},
    ylabel={$E_n/E_1$},
    every axis y label/.style={at={(axis description cs:-0.16,0.5)},
                               rotate=90, anchor=south, font=\small},
    xmin=0.4, xmax=5.6, ymin=0, ymax=26.5,
    xtick={1,2,3,4,5}, ytick={0,5,10,15,20,25},
    grid=major,
    legend pos=north west,
    legend style={font=\scriptsize, draw=light-gray},
    title={(b) \tpcifenglish{Spectrum}{Spektrum}},
]
  % Kasten: E_n / E_1 = n^2
  \addplot[cgblue, line width=1.0pt, mark=*, mark size=1.9pt,
           domain=1:5, samples=5] {x^2};
  \addlegendentry{\tpcifenglish{box}{Kasten}, $n^2$}

  % Oscillator, normalised the same way: E_n/E_1 = (n + 1/2)/(3/2)
  \addplot[birrenorange, line width=1.0pt, densely dashed,
           mark=o, mark size=1.9pt, mark options={solid},
           domain=1:5, samples=5] {(2*x+1)/3};
  \addlegendentry{\tpcifenglish{oscillator}{Oszillator}, $\tfrac{2n+1}{3}$}
\end{axis}

\end{tikzpicture}
\end{center}


The figure shows both at once: on the left each eigenfunction sits on its own
energy level, and one reads off that $\phi_n$ has exactly $n-1$ zeros in the
interior; on the right is the spectrum, both sequences normalised to their
respective $E_1$. The levels spread apart as $n^2$ -- in contrast to the
harmonic oscillator, whose levels we shall later find to be
\emph{evenly spaced}. So the shape of the potential determines not only
\emph{that} there is quantisation, but also \emph{how}.

The $\phi_n$ form a basis of our Hilbert space $\mathcal{H} = L^2([0,L])$ and thereby define a unitary
operator $U:\mathcal{H} \rightarrow \ell^2(\mathbb{N})$, namely $(U\Psi)(n) = \langle \phi_n|\Psi\rangle$, which diagonalises $H$. The time-evolution operator $U_t = e^{-itH/\hbar}$ is therefore obtained by expanding in the basis $\phi_n$ (= applying $U$), applying $U_t$ trivially to the basis vectors (= multiplication by a function of $n$) and summing the series with the new components (= applying $U^\dag$):
\begin{equation}
	U_t |\Psi\rangle = \sum_{n=1}^{\infty} e^{-\frac{it}{\hbar} E_n} |\phi_n\rangle \langle \phi_n|\Psi\rangle. 
\end{equation}
So as a mnemonic in Dirac's bra-ket notation we have once again
\begin{equation}
	\eqbox{U_t = \sum_{n=1}^{\infty} e^{-\frac{it}{\hbar} E_n} |\phi_n\rangle \langle \phi_n|.}
\end{equation}
With that, all five postulates have been ticked off for this system. The energy expectation value $\tr(\rho\widehat{H})$ can never fall below $E_1$, the energy of the \emph{ground state} $|\phi_1\rangle$; in the eigenstate $|\phi_n\rangle$ the variance moreover vanishes, $\tr(\rho\widehat{H}^2) - (\tr(\rho\widehat{H}))^2 = E_n^2 - E_n^2 = 0$. In the laboratory these levels are almost never observed directly, but rather as \emph{differences}: through photons given off in transitions between eigenstates.

\begin{intuition}
	One often reads that the $\phi_n$ are ``the'' states of the system. That is an idealisation. An energy eigenstate has no dynamics: up to a phase the particle would stay in it forever. In nature, however, there are always interactions, first and foremost with the electromagnetic field; after a while one has a superposition of several eigenstates. The eigenvectors are a computational device, not a typical state.
\end{intuition}

For a three-dimensional particle in a box shaped like a rectangular cuboid, the
essential work has already been done. Let us therefore consider the region
\begin{equation}
	G = \left\{(x,y,z)\,\middle| \,  0\le x \le L_x, 0\le y \le L_y, 0\le z \le L_z  \right\},
\end{equation}
where $L_x$, $L_y$ and $L_z$ are the edge lengths of the cuboid. The Hamiltonian is
now 
\begin{equation}
	\widehat{H} = -\frac{\hbar^2}{2m}\nabla^2,
\end{equation}
again with the boundary condition that functions in the domain of $H$ vanish on the boundary of the region. To solve the eigenvalue equation one makes a \emph{separation ansatz}
\begin{equation}
	\Psi(x,y,z) = \Phi_x(x)\Phi_y(y)\Phi_z(z)
\end{equation}
and finds that the function $\Phi_x''/\Phi_x$ no longer depends on $x$ (similarly for $\Phi_y$ and $\Phi_z$), so that $\Phi_x$ must be an eigenfunction for the problem of a one-dimensional interval just discussed. In this way we obtain the eigenfunctions and eigenvalues
\begin{align}
	\phi_{n_x,n_y,n_z}(x,y,z) &= \sqrt{\frac{8}{L_xL_yL_z}} \sin\left(\frac{n_x\pi}{L_x} x\right)\sin\left(\frac{n_y\pi}{L_y} y\right)\sin\left(\frac{n_z\pi}{L_z} z\right) \\
	E_{n_x,n_y,n_z}&=\frac{\hbar^2\pi^2}{2m} \left(\frac{n_x^2}{L_x^2} + \frac{n_y^2}{L_y^2} + \frac{n_z^2}{L_z^2} \right).
\end{align}
What can one say about boxes of other shapes? A separation ansatz works only in the few cases where one can find coordinates in which the shape of the box is easy to describe and in which the
Laplace operator moreover splits into a sum of second derivatives. Besides the cuboid one can treat the sphere in this way, and ellipsoids too. For general regions $G$ the determination of the eigenfunctions
and eigenvalues of the Laplace operator is not possible in closed form, but it is a problem that plays a role in various wave phenomena in physics (cavity resonance).

Let us take note: here, for once, a system has been worked through completely. That this succeeds is more a matter of luck than of rule -- the harmonic oscillator and the hydrogen atom are still to come, and then it stops.